\chapter{Justificación}

En el ámbito de los vehículos aéreos no tripulados de relativo bajo coste, como drones de uso personal, los controladores como Betaflight o iNav ya han resuelto el problema de la fusión de sensores en hardware de bajo coste. En estos sistemas se usan sensores y  microcontroladores muy populares y algoritmos relativamente fáciles de implementar en software. El hardware y software se encuentran en un estado muy maduro. El problema es que para uso educativo estos sistemas utilizan un software con millones de líneas de código que no los hacen intuitivos para estudiantes. Además, el PFD en estos sistemas no sigue las regulaciones presentes en normativa, cosa que tampoco los hace deseables para uso didáctico. \newline

Por otro lado, en las clases de ingeniería aeroespacial los estudiantes estudian la instrumentación de cabina de aeronaves (indicadores de actitud, altímetros...), a nivel teórico o mediante simuladores. Estos simuladores permiten familiarizarse con la interfaz de un PFD, pero detrás no hay hardware real, cosa que los hace menos intuitivos y más difícil de aprender de ellos. \newline

De estos dos factores nace la idea de este proyecto, una plataforma que conecte ambos mundos: usar el mismo tipo de hardware y algoritmos del ecosistema de drones para alimentar un PFD normalizado según la regulación aeronáutica. De esta forma, el estudiante puede recorrer el sistema completo de la instrumentación aviónica, desde la lectura del sensor hasta la representación en pantalla, incluyendo el hardware físico y el software. Además, el sistema completo se puede replicar por menos de 100\,€ en materiales, haciéndolo fácil de incorporar en laboratorios de forma sencilla, y todo el software empleado es de código abierto, lo que permite que otros estudiantes o docentes puedan reproducirlo sin coste de licencias. \newline

Como limitación principal, el sistema no está diseñado ni para certificación aeronáutica ni para el control real de un vehículo aéreo. Por este motivo no incorpora redundancia, protección ambiental ni los sistemas de seguridad que serían requeridos para un sistema real. Esto es coherente con su propósito: se trata de una plataforma de aprendizaje, no de un equipo de vuelo.